\clearpage
\VerdierPhysicalPageAnchor{0057}%
\VerdierSemanticAnchor{ch1-1-5-7-continuation}
五つの縦の図式を示している。それらは順に

\VerdierSemanticAnchor{ch1-1-5-7-vertical}
\[
  T'(\sigma')\Delta(\sigma,\tau),\quad
  T'(\sigma)\Delta(\sigma',\tau),\quad
  \Delta(\sigma,\tau)T(\sigma'),\quad
  \Delta(\sigma',\tau)T(\sigma),\quad
  \Delta(\sigma+\sigma',\tau)
\]
である。同様に四つの横の図式を示しており、それらは順に

\VerdierSemanticAnchor{ch1-1-5-7-horizontal}
\[
  D(\sigma,\sigma')T(\tau),\quad
  D(\sigma',\sigma)T(\tau),\quad
  T'(\tau)D(\sigma,\sigma'),\quad
  T'(\tau)D(\sigma',\sigma)
\]
である。

\VerdierSemanticAnchor{ch1-1-5-7-conclusion}
帰納法の仮定により、図式 $D(\sigma,\sigma')$ と
$D(\sigma',\sigma)$ はそれぞれ $\alpha(\sigma,\sigma')$-可換、
$\alpha(\sigma',\sigma)$-可換である。三つの場合を考えれば、主張
(1.5.4.3) の前半が従う。後半は $\sigma$ と $\tau$ の役割を入れ替えて
得られる。一般の場合に導入すべき図式はさらに複雑だが、恒等式
(1.2.1.3) を用いれば、論証は上で与えたものに帰着する。

\VerdierSemanticAnchor{ch1-1-5-8}
\noindent\textbf{注意 1.5.8.}\quad
命題 (1.5.2) は特に余サイクル $\alpha$ が自明な場合に適用できる。
このとき、圏 $\mathcal C'$ が加法圏でなくても、命題とその証明が
逐語的に拡張されることが分かる。

\VerdierSemanticAnchor{ch1-1-5-9}
\noindent\textbf{1.5.9.}\quad

\VerdierSemanticAnchor{ch1-1-5-9-functors}
\[
  (F_1,p_1),(F_2,p_2)\colon
  (\mathcal C,\mathbb Z^n,T,c)
  \longrightarrow
  (\mathcal C',\mathbb Z^n,T',c')
\]
を、それぞれ正規化された余サイクル
$\alpha_1,\alpha_2\in Z^2(\mathbb Z^n,\mathbb Z/2\mathbb Z)$
で捩られた二つの $\mathbb Z^n$-関手とする。また
$\beta\colon\mathbb Z^n\to\mathbb Z/2\mathbb Z$ を

\VerdierSemanticAnchor{ch1-1-5-9-eq1}
\begin{equation}
  \alpha_2(\sigma_1,\sigma_2)\alpha_1^{-1}(\sigma_1,\sigma_2)
  =\beta(\sigma_2)\beta(\sigma_1\sigma_2)^{-1}\beta(\sigma_1),
  \qquad \sigma_1,\sigma_2\in\mathbb Z^n,
  \tag{1.5.9.1}
\end{equation}
を満たす写像とし、$u\colon(F_1,p_1)\to(F_2,p_2)$ を関手の
$\beta$-射 (1.4.6) とする。$e_1,\ldots,e_n$ を $\mathbb Z^n$ の標準基底と
する。各 $i\in[1,n]$ に対し、次の図式は $\beta(e_i)$-可換である：

\VerdierSemanticAnchor{ch1-1-5-9-eq2}
\begin{equation}
  \resizebox{0.78\textwidth}{!}{%
    $\begin{tikzcd}[ampersand replacement=\&,column sep=6em,row sep=4em]
      F_1T(e_i)
        \arrow[r,"{p_1(e_i)}"]
        \arrow[d,"{u*T(e_i)}"']
        \& T'(e_i)F_1
          \arrow[d,"{T'(e_i)*u}"] \\
      F_2T(e_i)
        \arrow[r,"{p_2(e_i)}"]
        \& T'(e_i)F_2
    \end{tikzcd}$%
  }\;.
  \tag{1.5.9.2}
\end{equation}
